% ============================================================================
\section{Proof of Theorem S (Lemma S in a regime)}\label{app:S}
% ============================================================================
\intshare{Adapted from project note \texttt{math/07}, which survived adversarial review (\texttt{math/01} \S9). }{}Model M1, $b(t_f)$ free, $T:=t_f$,
$\Omega=\omega T$, $\tau(s)=\delta+\omega s$, $\bar\tau=\delta+\Omega$, $a_{10}=\rho_1/m_0$, $a_2=\rho_2/(m_0-\alpha\rho_2T)$,
$\varepsilon_m=a_1(T)/a_{10}-1$. The lateral force is $F=a\,d_h$, and the \emph{gain} of a control is $\int a_1-I_z$, the vertical
impulse wasted relative to upright minimum throttle. $P$ is the $L^2((0,T);\R^2)$ projection onto affine functions. It depends
only on the moments $m(F)=(\int F,\int sF)$, and $\pi_m$ denotes the affine function with moments $m$. In particular
$\norm{\pi_{(0,q)}}^2=12|q|^2/T^3$.

\paragraph{Landing from vertical data.} From $\zeta=(h\ez,v_{0z}\ez)$ a trajectory lands at $(p,0)$ with $v(T)=0$ iff $I_h=0$,
$I_z=g_0T-v_{0z}$, $\int sF=-p$ and $J_z$ is fixed by $\zeta$. Hence \emph{every trajectory from $\zeta$ has the same gain}, and
$m(F)=(0,-p)$.

\begin{proof}[Proof of \cref{lem:W}]
Let $\tilde f=f-f(0)-\tfrac sT(f(T)-f(0))$. Then $f-Pf=\tilde f-P\tilde f$, and $\norm{\tilde f-P\tilde f}\le\norm{\tilde f-\bar{\tilde f}}$ for the mean
$\bar{\tilde f}$. Since $\tilde f(0)=\tilde f(T)=0$, $\tilde f$ extends periodically, and the periodic Wirtinger inequality (Parseval) gives
$\norm{\tilde f-\bar{\tilde f}}^2\le(T/2\pi)^2\norm{\dot{\tilde f}}^2=(T/2\pi)^2(\norm{\dot f}^2-|f(T)-f(0)|^2/T)$. All steps are componentwise.
\end{proof}

\paragraph{Master inequality.} With $c_0=1/(1+\cos\delta)$, $\delta_e=2\sin(\delta/2)+\varepsilon_m\sin\bar\tau$,
$E_4=a_{10}(1+\varepsilon_m)^2\int_0^T(\frac1{1+\cos\tau(s)}-c_0)\sin^2\tau(s)\,ds$, $B=a_{10}c_0T(\frac\Omega{2\pi}+\delta_e)^2+E_4$,
$A'=c_0/a_{10}$, $\kappa_0=\cos\bar\tau/((a_2-a_{10})\sin^2\bar\tau)$ and $C_m=A'\kappa_0/(\kappa_0-A')$ ($C_m=A'$ if $\rho_1=\rho_2$):
\begin{proposition}[Master inequality {\pnote{math/07, Prop.~M}}]\label{prop:master}
% CLAIMS: S-M (CONFIRMED by verify-imp; FORMALIZED for constant mass) | Lean: master_inequality (variant (c), constant mass), master_inequality_fixed, dent (Slice/Master.lean; PENDING)
If $\bar\tau<\pi/2$ and $\kappa_0>A'$, every admissible M1 trajectory (any 3-D attitude path, throttle, $\alpha\ge0$) satisfies
$\mathrm{gain}\le B+C_m\norm{\pi_{m(F)}}^2$.\Sproved\ \Sleanpart{constant mass, variant with constant $\bar c(T)$}
\end{proposition}
\begin{proof}[Proof sketch]
(1) A pointwise majorant $-a\,d_z\le-a_1+a_1\bar c|d_h|^2-(a-a_1)\cos\tau$ with $\bar c=1/(1+\cos\tau)$ splits off the throttle surplus
$K=(a-a_1)d_h$, which costs at least $\kappa_0\norm K^2$. (2) Normalize the mass. (3) Eliminate $K$ through its projection
($PK=\pi_m-a_{10}Pv$) by maximizing a concave quadratic. (4) The only nonconvex step bounds the non-affine part of the lateral tilt
$\beta=b_h$, which is $\omega$-Lipschitz, by \cref{lem:W}: $\norm{P^\perp\beta}\le\omega T^{3/2}/(2\pi)$. The gimbal and mass corrections are
bounded pointwise by $\delta_e$.
\end{proof}

\paragraph{Witness.} A planar single-crossing S-turn at minimum throttle: tilt $\varphi(s)=\omega s$ for $s\le t_1$ and $\omega(2t_1-s)$
after, the thrust deflected outward by $\delta$ ($\chi=\varphi+\delta\,\mathrm{sgn}\,\varphi$), and $t_1$ chosen by the intermediate value
theorem so that $\int a_1\sin\chi=0$ (so $I_h=0$). Its mean-square lateral tilt is $(3\sqrt2-4)/6\cdot\Omega^2\approx0.0404\,\Omega^2$ in the
small-angle limit ($t_1=(1-1/\sqrt2)T$). This exceeds the Wirtinger cap $\Omega^2/(4\pi^2)\approx0.0253\,\Omega^2$. The witness does not return to
zero tilt, which is why it beats every path whose two lateral moments vanish.

\begin{proof}[Proof of \cref{thm:S}]
(1) Along the witness $\dot v_z=a_1\cos\chi-g_0>0$ by (S1), so $v_z<0$ on $[0,T)$ and $r_z>0$ before $T$. Rotations about $\ez$ preserve
$\zeta_w$, so $L(\zeta_w)$ contains the circle $|p|=|p_w|$. (2) Every trajectory from $\zeta_w$ landing at $(p,0)$ has gain $\mathrm{gain}_w$
and $m(F)=(0,-p)$, so \cref{prop:master} gives $\mathrm{gain}_w\le B+12C_m|p|^2/T^3$, i.e.\ $|p|\ge r_S$. (3) Apply \cref{thm:D} with
$\varphi(x)=\norm{x_r}+\iota_{\{x_v=0\}}(x)$; since $0\in\conv L(\zeta_w)$, the relaxed value is $0$.
\end{proof}

\begin{table}[h]\centering\small
\caption{Certified instances of \cref{thm:S} for a demonstration vehicle ($m_0=100$~kg, $\rho_1=1200$~N, $\rho_2=2500$~N, $\alpha=0$)\intshare{,
from \texttt{math/07} \S4}{}.}\label{tab:S}
\begin{tabular}{@{}lccccc@{}}\toprule
$\delta$ & $T$ [s] & $\Omega$ [rad] & margin $M$ [m/s] & dent $r_S$ [m] & witness offset $|p_w|$ [m]\\\midrule
$0$ & 2 & 0.40 & $1.39\times10^{-2}$ & 0.45 & 0.94\\
$0$ & 4 & 0.40 & $2.78\times10^{-2}$ & 1.80 & 3.74\\
$5^\circ$ & 2 & 0.30 & $1.57\times10^{-2}$ & 0.48 & 1.79\\
$10^\circ$ & 2 & 0.20 & $1.06\times10^{-2}$ & 0.39 & 2.59\\
$10^\circ$ & 4 & 0.20 & $2.11\times10^{-2}$ & 1.57 & 10.38\\\bottomrule
\end{tabular}
\end{table}
For $\delta>0$ the certified range of $\Omega$ has a positive lower end: as $\Omega\to0$ at fixed $\delta$, $\mathrm{gain}_w\to a_{10}T(1-\cos\delta)$
while $B\to a_{10}T\cdot2\sin^2(\delta/2)\cdot2/(1+\cos\delta)$, which is larger, so $M<0$. At $T=2$~s the certified ranges are $(0,0.55]$
($\delta=0$), $[0.004,0.46]$ ($5^\circ$), $[0.009,0.42]$ ($7^\circ$), $[0.038,0.34]$ ($10^\circ$) and $[0.070,0.27]$ ($12^\circ$) (independent \intshare{referee}{re-}
computation). All rows of \cref{tab:S} lie inside them.
With $\alpha=5.66\times10^{-4}$ the $\delta=5^\circ$ row fails narrowly, because $\varepsilon_m$ enters $\delta_e$ linearly. That is a weakness of the constants,
not of the geometry. For $\delta\ge15^\circ$ the analytic constants do not certify; a discretized S-procedure SDP suggests $\delta=15^\circ$
holds with an 8\% margin (numerical).
